\begin{table}[htbp]
\centering
\caption{Narrative frame-decomposition estimates}
\label{tab:h2_frame_decomposition}
\begin{tabular}{lrrrrc}
\toprule
Frame & Prevalence & \(b\): Frame \(\rightarrow \Delta\) & SE & \(p\) & LOMO sign stable \\
\midrule
Access Barriers & 52.7% & 8.83 & 0.54 & 0.000 & Yes \\
Collective Responsibility & 59.2% & -6.20 & 0.54 & 0.000 & Yes \\
Coercive Backlash & 8.4% & 6.02 & 0.97 & 0.000 & Yes \\
\bottomrule
\end{tabular}
\begin{flushleft}
\footnotesize Notes: The table reports descriptive frame-outcome estimates using high-burden rationales in the expanded fifteen-model sample. The dependent variable is the profile-repetition-level burden effect \(\Delta = W^{low} - W^{high}\). Coefficients are from a model including PVC and the three retained frame indicators. Positive coefficients indicate larger predicted burden effects. The table emphasizes frame-outcome associations rather than formal mediation. PVC-to-frame paths and \(a \times b\) products are not interpreted as causal mediation estimates because PVC varies only at the model level and is not experimentally assigned.
\end{flushleft}
\end{table}
